\section{Audit of Gierer--Meinhardt-Type Models}
\label{sec:gm-model-audit}

% TODO: add references on classical Gierer-Meinhardt models.
% TODO: add references on Hydra activator-inhibitor models.
% TODO: check notation against the source-density section once the final model family is fixed.

This section records the model variants that will be compared in the paper. The
point is not to rederive the general theory of diffusion-driven instability, but
to identify which qualitative properties are already present in the underlying
Gierer--Meinhardt equations and which properties are introduced by source
density.

\subsection{Two-Equation Activator--Inhibitor Models}
\label{subsec:two-equation-ai-models}

We start from the broad two-equation Gierer--Meinhardt activator-inhibitor
family
\begin{modelEquation}[eq:general-dimensional-gm]
    \frac{\partial u}{\partial t}
    &=
    D_u\Delta u
    +
    \frac{a u^2+q_u}{v}
    -\mu_u u
    +p_u,
    \\
    \frac{\partial v}{\partial t}
    &=
    D_v\Delta v
    +b u^2
    -\mu_v v
    +p_v.
\end{modelEquation}
Here \(u\) denotes the activator, \(v\) denotes the inhibitor, and homogeneous
Neumann boundary conditions are imposed. The terms \(p_v\), \(p_u\), and \(q_u\)
play different formal roles:
\begin{itemize}
    \item \(p_v\) is basal inhibitor production;
    \item \(p_u\) is additive activator production;
    \item \(q_u\) is activator production inside the inhibitor-regulated term.
\end{itemize}
The qualitative behaviour of the model depends strongly on which of these terms
are present.

\paragraph{Basal inhibitor production as a denominator shift.}
The term \(p_v\) is equivalent to adding a positive constant to the inhibitor
denominator. Indeed, writing the original inhibitor variable as \(V\) and setting
\begin{modelEquation}
    V
    &=
    v+K,
    \\
    K
    &:=
    \frac{p_v}{\mu_v},
\end{modelEquation}
the system can be written equivalently as
\begin{modelEquation}
    \frac{\partial u}{\partial t}
    &=
    D_u\Delta u
    +
    \frac{a u^2+q_u}{K+v}
    -\mu_u u
    +p_u,
    \\
    \frac{\partial v}{\partial t}
    &=
    D_v\Delta v
    +b u^2
    -\mu_v v.
\end{modelEquation}
In this shifted variable, \(v=0\) corresponds to the basal inhibitor level
\(V=K\).

\subsubsection{\texorpdfstring{No production terms: \(p_v=p_u=q_u=0\)}{No production terms}}

If all production terms are absent, the model is
\begin{modelEquation}
    \partial_t u
    &=
    D_u\Delta u
    +
    \frac{a u^2}{v}
    -\mu_u u,
    \\
    \partial_t v
    &=
    D_v\Delta v
    +b u^2
    -\mu_v v.
\end{modelEquation}
The origin is not a regular resting state because the first equation is singular
at \(v=0\). The homogeneous steady-state equations give
\begin{modelEquation}
    v
    &=
    \frac{b}{\mu_v}u^2,
    \\
    0
    &=
    \frac{a u^2}{v}-\mu_u u,
\end{modelEquation}
and hence the unique positive homogeneous steady state
\begin{modelEquation}[eq:no-production-positive-state]
    \bar u
    &=
    \frac{a\mu_v}{b\mu_u},
    \\
    \bar v
    &=
    \frac{a^2\mu_v}{b\mu_u^2}.
\end{modelEquation}
At this state,
\begin{modelEquation}
    J(\bar u,\bar v)
    &=
    \begin{pmatrix}
        \mu_u
        &
        \displaystyle -\frac{a\bar u^2}{\bar v^2}
        \\
        2b\bar u
        &
        -\mu_v
    \end{pmatrix}.
\end{modelEquation}
Thus the kinetic trace is \(\mu_u-\mu_v\), and the kinetic determinant is
\(\mu_u\mu_v>0\). For the usual parameter regime \(\mu_u<\mu_v\), the positive
state is stable for the spatially homogeneous kinetics. Since the activator
self-feedback entry is positive, this branch is the standard pattern-forming
branch when the inhibitor diffuses sufficiently faster than the activator.

The important modelling point is that this version does not contain a stable
no-head/resting state. The model is built around a positive pattern-forming
state, so patterns are expected for broad classes of initial conditions. In
simulations this variant can also generate very large peaks or spike-like
profiles because inhibition is unsaturated.

\subsubsection{\texorpdfstring{Inhibitor production only: \(p_v>0\), \(p_u=q_u=0\)}{Inhibitor production only}}

Adding basal inhibitor production removes the singularity. By the change of
variables above, the system takes the saturated form
\begin{modelEquation}[eq:gm-inhibitor-production]
    \frac{\partial u}{\partial t}
    &=
    D_u\Delta u
    +
    \frac{a u^2}{K+v}
    -\mu_u u,
    \\
    \frac{\partial v}{\partial t}
    &=
    D_v\Delta v
    +b u^2
    -\mu_v v.
\end{modelEquation}
This will be the reference two-equation model in the later subsections. After
rescaling one may set \(K=1\), but we keep \(K\) visible because it records the
effect of inhibitor production.

The zero state \((0,0)\) is now a regular homogeneous steady state. Its Jacobian
is
\begin{modelEquation}
    J(0,0)
    &=
    \begin{pmatrix}
        -\mu_u & 0\\
        0 & -\mu_v
    \end{pmatrix},
\end{modelEquation}
so the zero state is stable, with or without diffusion.

For \(u>0\), the steady-state equations give
\begin{modelEquation}
    v
    &=
    \frac{b}{\mu_v}u^2,
    \\
    \frac{a u}{K+v}
    &=
    \mu_u.
\end{modelEquation}
Equivalently,
\begin{modelEquation}[eq:inhibitor-production-positive-polynomial]
    \frac{\mu_u b}{\mu_v}u^2-au+\mu_u K
    &=
    0.
\end{modelEquation}
Thus two positive homogeneous steady states exist if and only if
\begin{modelEquation}[eq:inhibitor-production-two-positive-condition]
    a^2
    &>
    \frac{4b\mu_u^2K}{\mu_v}.
\end{modelEquation}
They are
\begin{modelEquation}[eq:inhibitor-production-positive-states]
    u_\pm
    &=
    \frac{\mu_v}{2b\mu_u}
    \left(
        a
        \pm
        \sqrt{
            a^2-\frac{4b\mu_u^2K}{\mu_v}
        }
    \right),
    \\
    v_\pm
    &=
    \frac{b}{\mu_v}u_\pm^2.
\end{modelEquation}

At a positive steady state,
\begin{modelEquation}
    J(u_\pm,v_\pm)
    &=
    \begin{pmatrix}
        \mu_u
        &
        \displaystyle -\frac{\mu_u^2}{a}
        \\
        2bu_\pm
        &
        -\mu_v
    \end{pmatrix}.
\end{modelEquation}
The determinant is
\begin{modelEquation}
    \det J(u_\pm,v_\pm)
    &=
    \mu_u
    \left(
        \frac{2b\mu_u}{a}u_\pm-\mu_v
    \right).
\end{modelEquation}
The lower branch \(u_-\) has negative determinant and is therefore unstable. The
upper branch \(u_+\) has positive determinant; it is stable for the homogeneous
kinetics when \(\mu_u<\mu_v\). Since the activator self-feedback entry is
positive, this upper branch is the candidate Turing-unstable branch when the
inhibitor diffuses sufficiently fast.

This is the first variant that naturally supports a threshold-regeneration
interpretation:
\begin{itemize}
    \item small perturbations of the zero state decay;
    \item the lower positive state acts as a threshold;
    \item sufficiently strong localized perturbations can reach the upper
    pattern-forming branch.
\end{itemize}
In contrast to the unsaturated model, the denominator \(K+v\) also regularizes
the inhibition and helps keep solutions bounded.

\subsubsection{\texorpdfstring{Activator and inhibitor production: \(p_v>0\), \(p_u\geq0\), \(q_u\geq0\)}{Activator and inhibitor production}}

With inhibitor production already encoded by \(K>0\), the two-equation model with
activator production is
\begin{modelEquation}[eq:gm-activator-and-inhibitor-production]
    \frac{\partial u}{\partial t}
    &=
    D_u\Delta u
    +
    \frac{a u^2+q_u}{K+v}
    -\mu_u u
    +p_u,
    \\
    \frac{\partial v}{\partial t}
    &=
    D_v\Delta v
    +b u^2
    -\mu_v v.
\end{modelEquation}
If \(p_u>0\) or \(q_u>0\), then \((0,0)\) is no longer a steady state. The
homogeneous steady states satisfy
\begin{modelEquation}
    v
    &=
    \frac{b}{\mu_v}u^2
\end{modelEquation}
and
\begin{modelEquation}[eq:two-equation-production-polynomial]
    0
    &=
    -\frac{\mu_u b}{\mu_v}u^3
    +
    \left(
        a+\frac{p_u b}{\mu_v}
    \right)u^2
    -\mu_u K u
    +
    (q_u+Kp_u).
\end{modelEquation}
Thus the number of positive steady states depends on the production strength.

For small \(p_u\) and \(q_u\), the phase portrait is a perturbation of
\eqref{eq:gm-inhibitor-production}: the zero state is replaced by a small
positive resting state, and one still expects a threshold branch and an upper
pattern-forming branch. For sufficiently large activator production, the lower
resting branch and the threshold branch may disappear in a saddle-node
bifurcation. Then the only remaining homogeneous state is the positive
pattern-forming branch, and the system tends to form patterns without requiring a
wound-like perturbation.

The two activator production terms are not identical algebraically: \(q_u\)
enters the numerator, while \(p_u\) is added outside the inhibitor-regulated
term. However, near small \(u\) they both increase the effective activator
production. This explains why their effects can look similar in simulations even
though the steady-state polynomial is changed in different coefficients.

\subsection{Reference Two-Equation Model for Regeneration}
\label{subsec:reference-two-equation-model}

The reference two-equation model for the regeneration discussion is therefore
\eqref{eq:gm-inhibitor-production}, the inhibitor-production model without
activator production. It has:
\begin{itemize}
    \item a stable zero/no-head state;
    \item two positive steady states when
    \(a^2>4b\mu_u^2K/\mu_v\);
    \item a lower positive branch that is unstable;
    \item an upper positive branch that is kinetically stable when
    \(\mu_u<\mu_v\) and can become unstable in the presence of diffusion.
\end{itemize}
This model already contains the basic threshold logic needed for wound-induced
regeneration. Source density is therefore not needed for the existence of a
stable resting state or a threshold to activation.

\subsection{Three-Equation Model with Source Density}
\label{subsec:three-equation-source-density-model}

We now add a source-density variable \(s\). In the shifted inhibitor variable,
the source-density model with both activator production mechanisms is
\begin{modelEquation}[eq:sd-general-production-model]
    \frac{\partial u}{\partial t}
    &=
    D_u\Delta u
    +
    s\frac{a u^2+q_u}{K+v}
    -\mu_u u
    +p_u,
    \\
    \frac{\partial v}{\partial t}
    &=
    D_v\Delta v
    +b s u^2
    -\mu_v v,
    \\
    \tau\frac{\partial s}{\partial t}
    &=
    D_s\Delta s
    +u
    -s.
\end{modelEquation}
Here \(K>0\) again represents basal inhibitor production. The term \(q_u\) is
source-dependent because it is multiplied by \(s\), whereas \(p_u\) is an
external activator production term.

\subsubsection{\texorpdfstring{No activator production: \(p_u=q_u=0\)}{No activator production}}

When \(p_u=q_u=0\), the model is
\begin{modelEquation}[eq:sd-no-activator-production]
    \partial_t u
    &=
    D_u\Delta u
    +
    s\frac{a u^2}{K+v}
    -\mu_u u,
    \\
    \partial_t v
    &=
    D_v\Delta v
    +b s u^2
    -\mu_v v,
    \\
    \tau\partial_t s
    &=
    D_s\Delta s
    +u
    -s.
\end{modelEquation}
The zero state \((0,0,0)\) is a homogeneous steady state. The reaction Jacobian
at zero is triangular up to the \(u\)-to-\(s\) coupling:
\begin{modelEquation}
    J(0,0,0)
    &=
    \begin{pmatrix}
        -\mu_u & 0 & 0\\
        0 & -\mu_v & 0\\
        1/\tau & 0 & -1/\tau
    \end{pmatrix},
\end{modelEquation}
so the zero state is stable.

For a positive homogeneous steady state, the third equation gives \(s=u\), and
the second equation gives
\begin{modelEquation}
    v
    &=
    \frac{b}{\mu_v}u^3.
\end{modelEquation}
The first equation then yields
\begin{modelEquation}[eq:sd-no-production-positive-polynomial]
    \frac{\mu_u b}{\mu_v}u^3
    -au^2
    +\mu_u K
    &=
    0.
\end{modelEquation}
Thus source density changes the positive steady-state equation from the
quadratic equation \eqref{eq:inhibitor-production-positive-polynomial} to a
cubic equation. Two positive steady states exist when the minimum of the cubic is
negative, equivalently
\begin{modelEquation}[eq:sd-no-production-two-positive-condition]
    \mu_u K
    &<
    \frac{4a^3\mu_v^2}{27\mu_u^2b^2}.
\end{modelEquation}

At a positive steady state, the reaction Jacobian in the variable order
\((u,v,s)\) is
\begin{modelEquation}[eq:sd-positive-jacobian]
    J(u,v,s)
    &=
    \begin{pmatrix}
        \mu_u
        &
        \displaystyle -\frac{\mu_u^2}{a u}
        &
        \mu_u
        \\
        2bu^2
        &
        -\mu_v
        &
        bu^2
        \\
        1/\tau
        &
        0
        &
        -1/\tau
    \end{pmatrix},
    \\
    s
    &=
    u,
    \qquad
    v
    =
    \frac{b}{\mu_v}u^3.
\end{modelEquation}
The characteristic polynomial has the form
\begin{modelEquation}
    \chi(\lambda)
    &=
    \lambda^3+\alpha_1\lambda^2+\alpha_2\lambda+\alpha_3,
\end{modelEquation}
where
\begin{modelEquation}
    \alpha_1
    &=
    \mu_v+\frac{1}{\tau}-\mu_u,
    \\
    \alpha_2
    &=
    \frac{\mu_v-2\mu_u}{\tau}
    -\mu_u\mu_v
    +
    \frac{2b\mu_u^2}{a}u,
    \\
    \alpha_3
    &=
    \frac{\mu_u}{\tau}
    \left(
        \frac{3b\mu_u}{a}u-2\mu_v
    \right).
\end{modelEquation}
The lower positive branch lies below the critical value
\begin{modelEquation}
    u_c
    &=
    \frac{2a\mu_v}{3b\mu_u},
\end{modelEquation}
so \(\alpha_3<0\) and that branch is unstable. The upper positive branch lies
above \(u_c\), so it is the only candidate for kinetic stability. It is stable
for the homogeneous kinetics precisely when the Routh--Hurwitz inequalities
\begin{modelEquation}
    \alpha_1
    &>
    0,
    \\
    \alpha_2
    &>
    0,
    \\
    \alpha_3
    &>
    0,
    \\
    \alpha_1\alpha_2
    &>
    \alpha_3
\end{modelEquation}
hold. With suitable time scale and degradation parameters, this upper branch is
the source-density analogue of the upper pattern-forming branch of the
two-equation reference model \eqref{eq:gm-inhibitor-production}.

\subsubsection{\texorpdfstring{Source-dependent activator production: \(q_u>0\), \(p_u=0\)}{Source-dependent activator production}}

If \(q_u>0\) but \(p_u=0\), the zero state remains a steady state because the
extra activator production is multiplied by \(s\). Its \(u,s\) linearization is
\begin{modelEquation}
    \begin{pmatrix}
        -\mu_u & q_u/K\\
        1/\tau & -1/\tau
    \end{pmatrix}.
\end{modelEquation}
Thus the zero state remains stable when
\begin{modelEquation}[eq:sd-qu-zero-stability]
    q_u
    &<
    \mu_u K,
\end{modelEquation}
and loses stability when \(q_u\) is too large.

For positive homogeneous states, \(s=u\), \(v=b u^3/\mu_v\), and
\begin{modelEquation}[eq:sd-qu-positive-polynomial]
    \frac{\mu_u b}{\mu_v}u^3-au^2+\mu_u K-q_u
    &=
    0.
\end{modelEquation}
Small \(q_u\) lowers the activation threshold but can preserve the resting
state. Large \(q_u\) destroys the threshold-regeneration interpretation because
the zero state is no longer stable.

\subsubsection{\texorpdfstring{External activator production: \(p_u>0\)}{External activator production}}

If \(p_u>0\), then \((0,0,0)\) is no longer a steady state. Positive homogeneous
states satisfy \(s=u\), \(v=b u^3/\mu_v\), and
\begin{modelEquation}[eq:sd-pu-positive-equation]
    0
    &=
    u(a u^2+q_u)
    +
    \left(
        p_u-\mu_u u
    \right)
    \left(
        K+\frac{b}{\mu_v}u^3
    \right).
\end{modelEquation}
This is generally a quartic equation in \(u\). For small \(p_u\), the behaviour
is expected to remain close to the threshold case: the zero state is replaced by
a small positive resting state, and large perturbations may reach the upper
pattern-forming branch. For sufficiently large \(p_u\), the lower resting branch
and the threshold branch may disappear, leaving only a positive pattern-forming
state. In that regime the model produces patterns without requiring a wound-like
trigger.

\subsection{Summary of the Audit}
\label{subsec:model-audit-summary}

\begin{table}[htbp]
    \centering
    \caption{Qualitative audit of the model variants used in this paper.}
    \label{tab:model-audit-summary}
    \resizebox{\textwidth}{!}{%
        \begin{tabular}{llllp{0.38\textwidth}}
            \toprule
            Model & Resting state & Positive branches & Unstable branch & Interpretation \\
            \midrule
            Two equations, no production
            & No regular zero
            & One positive branch
            & None
            & Pattern-forming model without a no-head state; can produce large spikes. \\
            Two equations, inhibitor production only
            & Stable zero
            & Two positive branches
            & Lower positive branch
            & Threshold-regeneration mechanism; upper branch is the pattern-forming candidate. \\
            Two equations, small activator production
            & Small positive rest state
            & Typically three positive states
            & Middle branch
            & Threshold behaviour persists approximately. \\
            Two equations, large activator production
            & No resting/threshold regime
            & Typically one positive state
            & None
            & Patterning may occur without a wound-like perturbation. \\
            Source density, no activator production
            & Stable zero
            & Two positive branches when \eqref{eq:sd-no-production-two-positive-condition} holds
            & Lower positive branch
            & Similar threshold logic, but the steady-state equation and stability conditions change. \\
            Source density, source-dependent \(q_u\)
            & Stable zero if \(q_u<\mu_u K\)
            & Production-dependent
            & Production-dependent
            & Small \(q_u\) lowers the threshold; large \(q_u\) can destabilize the resting state. \\
            Source density, external \(p_u\)
            & Zero absent
            & Production-dependent
            & Production-dependent
            & Small \(p_u\) gives a small positive rest state; large \(p_u\) removes the threshold regime. \\
            \bottomrule
        \end{tabular}%
    }
\end{table}

The conclusion of this audit is that the inhibitor-production version of the
two-equation Gierer--Meinhardt model already gives a stable resting state, a
threshold to activation, and a pattern-forming upper branch. Source density does
not need to be introduced to obtain basic regeneration-like threshold dynamics.
Its role should therefore be tested in the later sections as a mechanism for
positional memory, grafting response, and long-lived anchoring of head-forming
regions.
